\documentclass{article}
\usepackage{listings}
\usepackage{amsmath,amsthm}
\usepackage{amssymb}
\usepackage{hyperref}
\usepackage{xurl}


\title{Group projects Programming in Python II}
\author{}
\date{}

\begin{document}
\maketitle

\textbf{Deadline: 14/12/2024, 23:59,  please submit your *.py or *.ipynb files using the designated folder on the K2 site of the course.}



\section*{Project 1: Group 6 (Maxime Yilmaz, Alexandre Paquet, Christian Cheng), Group 7 (Zuer Li,Qiuyan Li,LUONG Thi Lan Hue) }
The task is to predict stock prices based on past stock prices. The data is stored in the file \lstinline$tech_companies_stock_prices.csv$.

The fields are as follows:
\begin{itemize} \item \lstinline$Date$ -- the date of the transaction 
\item \lstinline$Stock_name$ -- the name of the company 
\item \lstinline$Open$, \lstinline$Close$ -- opening and closing prices on the transaction date 
\item \lstinline$Open_i$, \lstinline$Close_i$ -- for $i=1,2,3,4,5$, these represent the opening and closing prices on \lstinline$Date$-$i$ 
\end{itemize}

For each company, predict the share price based on its past share prices over the previous $k$ dates. Specifically: 
\begin{itemize} 
\item Try to predict both the opening and closing prices using various choices of past opening and closing prices. 
\item Explore using past share prices of other companies to predict the share price of a given company. 
\end{itemize}

Use various models (such as linear regression, SVM, decision trees, and neural networks) to predict share prices. Compare the models in terms of prediction accuracy to determine which model performs best.

Finally, select one or two independent variables and use them to build models for predicting share prices. Identify which variables provide the best model for predicting share prices.


%% The task is to predict stock prices based on the past stock prices. The .csv is \lstinline$tech_companies_stock_prices.csv$. 
%%  The fields are as follows 
%%  \begin{itemize}
%%    \item \lstinline$Date$ -- date of transaction
%%    \item \lstinline$Stock name$ -- name of the company
%%    \item \lstinline$Open$, \lstinline$Close$ -- opening, closing prices.
%%    \item \lstinline$Open_i$, \lstinline$Close_i$
%%	    $i=1,2,3,4, 5$ opening, closing prices at \lstinline$Date$-$i$. 
%%  \end{itemize}
%%
%%   For each company, predict the share price for this company based on the past share price of the same company during the previous $k$ dates. 
%%   Try to predict  opening, closin using various choices of past opening  prices. 
%%   Try to use the past share price of other companies to predict the share price of a given company. 
%%
%%    Use various models (linear regression, SVM, decision tress, neural networks) for predicting the share price.
%%    Compare the various models in terms of accuracy of prediction. Which one is the best ? 
%%    Select one or two independent variables and use them to build models for predicting share prices. Which variables
%%    give the best model for predicting share prices ?




\section*{Project 2: Group 8 (Killian Bruni,Maïlys Texier), Group 9 (ZHANG Ruge,CAO Ranjie) }
The task is to predict stock prices based on past stock prices. The data is stored in the file \lstinline$stock_data_for_hw1.csv$.

The fields are as follows: \begin{itemize} \item \lstinline$Date$ -- date of the transaction \item \lstinline$Company$ -- name of the company \item \lstinline$Open$, \lstinline$Close$, \lstinline$High$, \lstinline$Low$ -- opening, closing, highest, and lowest prices on the transaction date \item \lstinline$Open_i$, \lstinline$Close_i$, \lstinline$High_i$, \lstinline$Low_i$ -- for $i=1,2,3,4,5$, these represent the opening, closing, highest, and lowest prices on \lstinline$Date$-$i$ \end{itemize}

For each company, predict the share price based on its past share prices over the previous $k$ dates. Specifically: \begin{itemize} \item Try to predict opening, closing, high, and low prices separately, using different combinations of past opening, closing, high, and low prices. \item Investigate using past share prices of other companies to predict the share price of a given company. \end{itemize}

Use various models (e.g., linear regression, SVM, decision trees, neural networks) for predicting share prices. Compare the models in terms of prediction accuracy to determine which one performs best.

Finally, select one or two independent variables to build models for predicting share prices. Identify which variables yield the best predictive model.

\section*{Project 3: Group 1 (Anouk Tourtelier,Lana Szymoniak,Florine Rousteau), Group 3 (Vincent MALRANGE,Victor Bouffant,Léo Rioufrays--Laplace) }
 Predict stock prices based on past stock prices, ESG Score, and net income. The CSV file \lstinline$StockESG0.csv$ contains the following columns: \begin{itemize} \item \lstinline$Date$ -- date of transaction \item \lstinline$Company$ -- name of the company \item \lstinline$Open$, \lstinline$Close$ -- opening and closing prices on the transaction date \item \lstinline$Open_i$, \lstinline$Close_i$ -- for $i=1,2,3,4,5$, these represent the opening and closing prices on \lstinline$Date$-$i$ \item \lstinline$Net income$ -- net income of the company \item \lstinline$ESG Score$ -- overall ESG score \item \lstinline$ESG Controversies Score$, \lstinline$Environmental Controversies Count$, \lstinline$Human Rights Controversies$, \lstinline$Business Ethics Controversies$, \lstinline$Wages Working Condition Controversies Count$, \lstinline$Total Renewable Energy To Energy Use in million$, \lstinline$Estimated CO2 Equivalents Emission Total$ -- various ESG-related scores and metrics \end{itemize}

Use past opening or closing prices, net income, and various ESG scores to predict the current opening (or closing) price.

Apply different models, such as linear regression, SVM, decision trees, and neural networks, for predicting the share price. Compare the models in terms of prediction accuracy to determine which performs best.

Finally, select one or two independent variables to build predictive models for share prices. Identify which variables yield the most accurate model. 



 \section*{Project 4:  Group 2 (Thelma GRECO,Angelick ORANGE), Group  5 (Yipan Liu,Ruixuan Li,Jingyan Wu) }
The task is to predict the ESG score based on various other variables, such as: \begin{itemize} \item \lstinline$ESG Controversies Score$ \item \lstinline$Environmental Controversies Count$ \item \lstinline$Human Rights Controversies$ \item \lstinline$Business Ethics Controversies$ \item \lstinline$Wages Working Condition Controversies Count$ \item \lstinline$Total Renewable Energy To Energy Use in million$ \item \lstinline$Estimated CO2 Equivalents Emission Total$ \end{itemize}

The data is stored in the file \lstinline$ESGAllCompanies1.csv$.

Use various models (e.g., linear regression, SVM, decision trees, neural networks) to predict the ESG score. Compare the performance of the models in terms of prediction accuracy. Which model performs the best?

Select one or two independent variables and build models to predict the ESG score. Identify which variables provide the best predictive model for the ESG score. 

\section*{Project 5: Group 5 (Célia Dompiétrini, Loane Mboutou,  Zaineb Er-raji)}
The task is to predict the logarithm of house prices (variable \emph{lprice}) using the following variables: \begin{itemize} \item logarithm of the lot size (variable \emph{llotsize}) \item logarithm of the surface area of the house (variable \emph{lsqrft}) \item number of bedrooms (variable \emph{bdrms}) \item surface area of the house (variable \emph{sqrft}) \item lot size (variable \emph{lotsize}) \end{itemize}

The data is stored in the file \lstinline$hprice1_merged.csv$.

Use various models (such as linear regression, SVM, decision trees, and neural networks) to predict the logarithm of the house price. Compare the performance of the models in terms of prediction accuracy. Which model performs the best?

Select one or two independent variables and build models to predict the logarithm of the price. Identify which variables provide the best model for predicting the logarithm of the price.

Repeat the exercise by replacing the logarithm of the price with the price itself (variable \emph{price}).


\end{document}
